\chapter[Taban Aritmetiği · \textit{İleri}]{Taban Aritmetiği}
\chaptermark{Taban Aritmetiği}

Günlük hayatta sayıları neredeyse istisnasız onluk sistemde yazar ve düşünürüz: $354$ sayısı bizim için üç yüz elli dörttür. Ancak doğanın veya matematiğin kendisinde $10$ sayısına tanınmış nesnel bir ayrıcalık yoktur. İnsanlığın onluk sistemi benimsemesindeki temel etken, iki elimizde toplam on parmağımızın bulunmasıdır. Şayet sekiz parmağımız olsaydı, bugün muhtemelen sekizlik sistemi kullanıyor olacaktık. 

Bilgisayar biliminde ise fiziksel gerçeklik farklıdır: Sayısal devreler transistörlerden oluşur ve bir transistör temel olarak iki kararlı duruma sahiptir: ``elektrik akımı var'' (yüksek voltaj, $1$) veya ``elektrik akımı yok'' (düşük voltaj, $0$). Bu nedenle modern bilgisayarların ana dili ikili (binary) tabandır. Dahası, uzun ikili dizileri daha okunabilir kılmak adına sekizli (octal) ve onaltılı (hexadecimal) tabanlar tercih edilir.

Bu bölümde pozisyonel sayı sistemlerinin matematiksel temellerini, tabanlar arası dönüşüm yöntemlerini, farklı tabanlarda aritmetik işlemleri ve ikili sistemin algoritmik problem çözmedeki yerini inceleyeceğiz.

\section{Taban ve Dönüşüm}

\subsection{Pozisyonel Sayı Sistemi ve Sezgi}

Bir sayıyı yazarken her bir rakamın değeri, bulunduğu basamağa (pozisyona) bağlıdır. Örneğin $10$ tabanında yazılmış $354$ sayısını basamak çözümlemesiyle ele alalım:
\[
354 = 3 \cdot 100 + 5 \cdot 10 + 4 \cdot 1 = 3 \cdot 10^2 + 5 \cdot 10^1 + 4 \cdot 10^0
\]
Buradaki temel ilke şudur: Her basamak, tabanın ($10$'un) artan bir kuvvetini temsil eder ve kullanılan rakamlar daima $\{0, 1, 2, \dots, 9\}$ kümesindendir; yani tabandan kesinlikle küçüktür.

Tabanımız $10$ yerine $2$ olsaydı basamak değerleri $2$'nin kuvvetleri ($2^0=1, 2^1=2, 2^2=4, 2^3=8, \dots$) olur ve kullanabileceğimiz rakamlar yalnızca $\{0, 1\}$ kümesinden seçilirdi. Örneğin $(101)_2$ sayısını açarsak:
\[
(101)_2 = 1 \cdot 2^2 + 0 \cdot 2^1 + 1 \cdot 2^0 = 4 + 0 + 1 = 5
\]
elde ederiz. Benzer biçimde taban $7$ seçilseydi basamak ağırlıkları $7^0, 7^1, 7^2, \dots$ olacak ve izin verilen rakamlar $\{0, 1, 2, 3, 4, 5, 6\}$ kümesini oluşturacaktı.

Bu sezgiyi matematiksel bir tanıma dönüştürelim.

\begin{definition}[Sayı Tabanı ve Pozisyonel Gösterim]
$b \ge 2$ bir tamsayı olsun. Bir $N$ doğal sayısının $b$ tabanındaki (radix-$b$) gösterimi, $a_k \ne 0$ ve her $0 \le i \le k$ için $a_i \in \{0, 1, \dots, b-1\}$ olmak üzere,
\[
(a_k a_{k-1} \dots a_1 a_0)_b
\]
biçiminde ifade edilen dizidir. Bu gösterimin sayısal değeri
\[
N = \sum_{i=0}^{k} a_i b^i = a_k b^k + a_{k-1} b^{k-1} + \dots + a_1 b + a_0
\]
olarak tanımlanır. Taban belirtilmemişse varsayılan taban $10$'dur.
\end{definition}

Taban $10$'dan büyük olduğunda ($b > 10$), $9$'dan büyük rakamları tek bir sembolle göstermek için harfler kullanılır: $A=10, B=11, C=12, D=13, E=14, F=15, \dots$ gibi. Örneğin onaltılı tabanda $(2F)_{16} = 2 \cdot 16^1 + 15 \cdot 16^0 = 32 + 15 = 47$ olur.

Her pozitif tamsayının, verilen herhangi bir $b \ge 2$ tabanında tek bir şekilde yazılabileceğini aşağıdaki teorem garanti eder.

\begin{theorem}[Taban Gösteriminin Varlığı ve Tekliği]
$b \ge 2$ sabit bir tamsayı olsun. Her pozitif $N$ tamsayısı için,
\[
N = a_k b^k + a_{k-1} b^{k-1} + \dots + a_1 b + a_0
\]
eşitliğini sağlayan, $a_k \ne 0$ ve her $i \in \{0, 1, \dots, k\}$ için $0 \le a_i < b$ koşuluna uyan bir $k \ge 0$ tamsayısı ve katsayılar dizisi $(a_k, a_{k-1}, \dots, a_0)$ tek bir biçimde mevcuttur.
\end{theorem}

\begin{proof}
İspatı iki aşamada kuralım: önce varlık (algoritmik inşa), ardından teklik.

\textbf{1. Varlık:} Bölüm 11'de gördüğümüz Bölme Algoritmasını ($N = q \cdot b + r$, $0 \le r < b$) ardışık olarak uygulayalım:
\begin{align*}
N &= q_0 b + a_0, \quad 0 \le a_0 < b \\
q_0 &= q_1 b + a_1, \quad 0 \le a_1 < b \\
q_1 &= q_2 b + a_2, \quad 0 \le a_2 < b \\
&\;\;\vdots \\
q_{k-1} &= q_k b + a_k, \quad 0 \le a_k < b \quad (q_k = 0 \text{ olduğunda durulur})
\end{align*}
$b \ge 2$ olduğundan, $N > q_0 > q_1 > q_2 > \dots \ge 0$ kesin azalan bir negatif olmayan tamsayılar dizisidir. İyisıralama ilkesi gereği bu dizi sonsuza kadar devam edemez; sonlu bir adımda bir $q_k = 0$ değerine ulaşmak zorundadır ($q_{k-1} \ne 0$ olduğundan $a_k \ne 0$'dır).

Elde edilen kalanları geriye doğru yerine yazarsak:
\begin{multline*}
N = q_0 b + a_0 = (q_1 b + a_1) b + a_0 = q_1 b^2 + a_1 b + a_0 \\
= \dots = a_k b^k + a_{k-1} b^{k-1} + \dots + a_1 b + a_0
\end{multline*}
ifadesine ulaşırız. Bu, aranan gösterimin var olduğunu gösterir.

\textbf{2. Teklik:} Gösterimin tek olmadığını varsayalım. Bu durumda $N$ sayısının $b$ tabanında iki farklı gösterimi bulunsun. Gerekirse katsayıların başına sıfır ekleyerek basamak sayılarını eşitleyebiliriz:
\[
N = \sum_{i=0}^{m} a_i b^i = \sum_{i=0}^{m} c_i b^i, \quad 0 \le a_i, c_i < b
\]
Buradan:
\[
\sum_{i=0}^{m} (a_i - c_i) b^i = 0
\]
yazılır. Gösterimler farklı olduğundan, $a_i \ne c_i$ şartını sağlayan en az bir indeks vardır. Bu indekslerin en küçüğüne $j$ diyelim ($0 \le j \le m$). O halde her $i < j$ için $a_i = c_i$'dir. Eşitlik şu biçimi alır:
\[
(a_j - c_j) b^j + \sum_{i=j+1}^{m} (a_i - c_i) b^i = 0 \implies (c_j - a_j) b^j = \sum_{i=j+1}^{m} (a_i - c_i) b^i
\]
Her iki tarafı $b^{j+1}$ modülünde incelersek, sağ taraf $b^{j+1}$'in tam katı olduğundan:
\[
(c_j - a_j) b^j \equiv 0 \pmod{b^{j+1}} \implies c_j - a_j \equiv 0 \pmod b
\]
elde edilir. Fakat $0 \le a_j, c_j < b$ olduğundan aralarındaki fark $-b < c_j - a_j < b$ aralığındadır. Bu aralıkta $b$'ye bölünebilen tek tamsayı $0$'dır. Demek ki $c_j - a_j = 0$, yani $a_j = c_j$ olmalıdır. Bu ise $j$'nin $a_j \ne c_j$ şartını sağlayan en küçük indeks olmasıyla çelişir. Dolayısıyla gösterim tek olmak zorundadır.
\end{proof}

\subsection{Tabanlar Arası Dönüşüm Yöntemleri}

\subsubsection{Keyfi bir tabandan 10 tabanına geçiş}

$(a_k a_{k-1} \dots a_1 a_0)_b$ sayısını $10$ tabanına dönüştürmek için tanım gereği polinom açılımı hesaplanır:
\[
N = a_k b^k + a_{k-1} b^{k-1} + \dots + a_1 b + a_0
\]
Bu toplamı doğrudan kuvvet alarak hesaplamak yerine, algoritmalarda sıkça başvurulan \textbf{Horner Yöntemi} ile paranteze alarak iç içe hesaplamak hem elle işlem yaparken hem de kod yazarken işlem sayısını ve hata payını azaltır:
\[
N = \Big( \dots \big((a_k \cdot b + a_{k-1}) \cdot b + a_{k-2}\big) \cdot b + \dots + a_1 \Big) \cdot b + a_0
\]

\begin{example}
$(110101)_2$ sayısını $10$ tabanına çeviriniz.
\end{example}

\begin{proof}[Çözüm]
\textbf{1. Yol (Kuvvetler toplamı):}
Basamak değerlerini sağdan sola doğru $2^0, 2^1, 2^2, 2^3, 2^4, 2^5$ olarak yazalım:
\[
(110101)_2 = 1 \cdot 2^5 + 1 \cdot 2^4 + 0 \cdot 2^3 + 1 \cdot 2^2 + 0 \cdot 2^1 + 1 \cdot 2^0
\]
\[
= 32 + 16 + 0 + 4 + 0 + 1 = 53
\]

\textbf{2. Yol (Horner yöntemi):}
En soldaki bitten başlayıp her adımda sonucu $2$ ile çarpıp bir sonraki biti ekleyelim:
\begin{align*}
\text{Başlangıç: } & 1 \\
\text{1. adım: } & 1 \cdot 2 + 1 = 3 \\
\text{2. adım: } & 3 \cdot 2 + 0 = 6 \\
\text{3. adım: } & 6 \cdot 2 + 1 = 13 \\
\text{4. adım: } & 13 \cdot 2 + 0 = 26 \\
\text{5. adım: } & 26 \cdot 2 + 1 = 53
\end{align*}
Her iki yöntem de beklediğimiz gibi $53$ sonucunu verir.
\end{proof}

\subsubsection{10 tabanından keyfi bir tabana geçiş: Ardışık Bölme}

Varlık ispatındaki kurgu, doğrudan pratik dönüşüm yöntemini verir:
Sayıyı $b$'ye böleriz; kalan, gösterimin en sağdaki basamağını ($a_0$) oluşturur. Elde edilen bölümü tekrar $b$'ye böldüğümüzde yeni kalan $a_1$ olur. Bu işleme bölüm $0$ olana kadar devam edilir. Kalanlar sondan başa (aşağıdan yukarıya) doğru dizilerek sayı elde edilir.

\begin{example}
$158$ sayısını $2$ ve $7$ tabanlarında yazınız.
\end{example}

\begin{proof}[Çözüm]
Önce $7$ tabanına dönüştürelim:
\begin{align*}
158 \div 7 &= 22 \quad (\text{kalan } a_0 = 4) \\
22 \div 7 &= 3 \quad (\text{kalan } a_1 = 1) \\
3 \div 7 &= 0 \quad (\text{kalan } a_2 = 3)
\end{align*}
Kalanları sondan başa doğru okursak: $158 = (314)_7$. Sağlamasını yapalım:
$3 \cdot 7^2 + 1 \cdot 7^1 + 4 \cdot 7^0 = 3 \cdot 49 + 7 + 4 = 147 + 11 = 158$.

Şimdi aynı sayıyı $2$ tabanına dönüştürelim:
\begin{align*}
158 \div 2 &= 79 \quad (\text{kalan } 0) \\
79 \div 2 &= 39 \quad (\text{kalan } 1) \\
39 \div 2 &= 19 \quad (\text{kalan } 1) \\
19 \div 2 &= 9 \quad (\text{kalan } 1) \\
9 \div 2 &= 4 \quad (\text{kalan } 1) \\
4 \div 2 &= 2 \quad (\text{kalan } 0) \\
2 \div 2 &= 1 \quad (\text{kalan } 0) \\
1 \div 2 &= 0 \quad (\text{kalan } 1)
\end{align*}
Kalanları aşağıdan yukarıya doğru yazarsak: $158 = (10011110)_2$.
\end{proof}

\subsubsection{$2^k$ Tabanları Arasındaki Doğrudan Geçiş (Gruplama İlkesi)}

Bilgisayarlar sayıları ikili tabanda (bitler hâlinde) saklayıp işlediğinden, bit düzeyindeki işlemler (\texttt{xor}, \texttt{and}, \texttt{or}, shift gibi) aritmetik işlemlere (mod, toplama, çarpma gibi) göre donanımda daha hızlı çalışır. Bu yakınlık, bilgisayar biliminde ikili ($2^1$), sekizli ($2^3$) ve onaltılı ($2^4$) sistemlerin yaygın olmasının da arkasındaki nedendir; üstelik bu tabanlar arasında $10$ tabanına uğramadan, basamakları bloklayarak doğrudan geçiş yapılabilir.

$b = 2^k$ olsun. İkili tabandaki bir sayının her $k$ basamaklık bloğu, doğrudan $2^k$ tabanındaki tek bir basamağa karşılık gelir. Çünkü:
\[
\sum_{i=0}^{k-1} a_i 2^i < \sum_{i=0}^{k-1} 1 \cdot 2^i = 2^k - 1
\]
olup, bu değer tam olarak $2^k$ tabanındaki bir rakamın alabileceği değer aralığındadır ($0$ ile $2^k-1$ arası).

\begin{itemize}
    \item \textbf{İkili $\leftrightarrow$ Sekizli ($2^3$):} İkili sayı sağdan sola doğru üçerli gruplara ayrılır (en soldaki eksik basamaklar sıfırla tamamlanır). Her $3$ bitlik grup bir sekizli rakama çevrilir.
    \item \textbf{İkili $\leftrightarrow$ Onaltılı ($2^4$):} İkili sayı sağdan sola doğru dörderli gruplara ayrılır. Her $4$ bitlik grup (nibble) onaltılı bir rakama ($0\text{--}9, A\text{--}F$) çevrilir.
\end{itemize}

\begin{example}
$(11010111101)_2$ sayısını sekizli ve onaltılı tabana çeviriniz.
\end{example}

\begin{proof}[Çözüm]
\textbf{Sekizli dönüşüm (üçerli gruplama):}
Sağdan sola üçer üçer ayıralım:
\[
(011 \mid 010 \mid 111 \mid 101)_2
\]
Her bir grubu sekizli rakama dönüştürelim:
\begin{align*}
(011)_2 &= 3 \\
(010)_2 &= 2 \\
(111)_2 &= 7 \\
(101)_2 &= 5
\end{align*}
Buradan sayı $(3275)_8$ bulunur.

\textbf{Onaltılı dönüşüm (dörderli gruplama):}
Sağdan sola dörder dörder ayıralım:
\[
(0110 \mid 1011 \mid 1101)_2
\]
Her bir grubu onaltılı rakama dönüştürelim:
\begin{align*}
(0110)_2 &= 6 \\
(1011)_2 &= 11 \implies B \\
(1101)_2 &= 13 \implies D
\end{align*}
Böylece $(6BD)_{16}$ elde edilir.
\end{proof}

\subsection{Basamak Sayısı ve Tabanın Büyüklüğü}

Bir $N$ pozitif tamsayısının $10$ tabanındaki basamak sayısını $\lfloor \log_{10} N \rfloor + 1$ formülüyle hesaplarız. Bu kural herhangi bir $b \ge 2$ tabanı için de geçerlidir.

\begin{theorem}[Basamak Sayısı]
$N$ pozitif bir tamsayı ve $b \ge 2$ olsun. $N$'nin $b$ tabanındaki basamak sayısı
\[
L_b(N) = \lfloor \log_b N \rfloor + 1
\]
ile verilir.
\end{theorem}

\begin{proof}
$N$ sayısı $b$ tabanında $k+1$ basamaklı olsun. Bu durumda en yüksek basamağın katsayısı $a_k \ge 1$ olmak üzere:
\[
b^k \le a_k b^k \le N = \sum_{i=0}^k a_i b^i \le \sum_{i=0}^k (b-1) b^i = (b-1) \frac{b^{k+1}-1}{b-1} = b^{k+1} - 1 < b^{k+1}
\]
Demek ki,
\[
b^k \le N < b^{k+1}
\]
Eşitsizliğin $b$ tabanında logaritmasını alırsak ($f(x) = \log_b x$ kesin artan fonksiyondur):
\[
k \le \log_b N < k + 1
\]
Tam değer (taban) tanımı gereği:
\[
k = \lfloor \log_b N \rfloor
\]
Basamak sayısı $k+1$ olduğundan, toplam basamak sayısı $\lfloor \log_b N \rfloor + 1$ olarak bulunur.
\end{proof}

Bu bağıntı, bir sayının küçük tabanlarda (örneğin $b=2$) daha çok, büyük tabanlarda (örneğin $b=16$) daha az basamakla temsil edildiğini gösterir. $2^{10} = 1024 \approx 10^3$ yaklaşımından hareketle, bir sayının ikili tabandaki basamak sayısı onluk tabandaki basamak sayısının yaklaşık $\log_2 10 \approx 3.32$ katıdır.

\subsection{Çözümlü Örnekler}

Olimpiyat sınavlarında taban aritmetiği genellikle cebirsel denklemler, basamak çözümlemeleri veya bölünebilme kurallarıyla bir arada sorulur.

\begin{example}
$(ab)_8 = (ba)_9$ eşitliğini sağlayan sıfırdan farklı $a$ ve $b$ rakamlarını bulunuz.
\end{example}

\begin{proof}[Çözüm]
İki farklı tabandaki bu ifadeyi ortak bir zemine çekmek için $10$ tabanındaki polinom açılımlarını kullanalım. Ayrıca taban tanımlarından gelen basamak sınırlarını ($a, b < 8$ ve $a, b < 9$) unutmamalıyız.

$(ab)_8$ sayısının $8$ tabanında iki basamaklı bir sayı belirtmesi için $1 \le a \le 7$ ve $0 \le b \le 7$ olmalıdır. Benzer şekilde $(ba)_9$ için $1 \le b \le 8$ ve $0 \le a \le 8$ gereklidir. Bu iki koşulun kesişimi:
\[
1 \le a \le 7 \quad \text{ve} \quad 1 \le b \le 7
\]
Açılımları yazarsak:
\[
(ab)_8 = 8a + b
\]
\[
(ba)_9 = 9b + a
\]
Eşitliği düzenleyelim:
\[
8a + b = 9b + a \implies 7a = 8b
\]
$\gcd(7, 8) = 1$ olduğundan (Bölüm 13), $7 \mid 8b \implies 7 \mid b$ olmalıdır.
Fakat $1 \le b \le 7$ aralığında $7$'nin tek katı $b = 7$'dir.
$b = 7$ yazılırsa $7a = 8 \cdot 7 \implies a = 8$ bulunur.

Ancak $a=8$ değeri $1 \le a \le 7$ aralığına düşmez; $8$ tabanında $8$ rakamı bulunamaz.
Dolayısıyla bu eşitliği sağlayan hiçbir $a, b$ rakam çifti \textbf{yoktur}.
\end{proof}

\begin{example}
Hangi $b > 4$ tabanında $(1331)_b$ sayısı bir tamküptür?
\end{example}

\begin{proof}[Çözüm]
$1331$ sayısındaki $1, 3, 3, 1$ katsayıları, Bölüm 6'dan hatırlayacağımız binom katsayıları $\binom{3}{k}$ ile örtüşmektedir.

Sayının $b$ tabanındaki açılımını yazalım:
\[
(1331)_b = 1 \cdot b^3 + 3 \cdot b^2 + 3 \cdot b + 1
\]
Bu cebirsel ifadeyi çarpanlarına ayırırsak:
\[
b^3 + 3b^2 + 3b + 1 = (b + 1)^3
\]
Görüldüğü üzere, $b > 3$ koşulunu sağlayan \textbf{bütün} tabanlarda $(1331)_b$ sayısı $(b+1)$ tamsayısının küpüne eşittir.
Dolayısıyla aranan tabanlar $b \ge 5$ (veya sorudaki kısıta göre $b > 4$) şartını sağlayan \textbf{tüm tamsayılardır}.
\end{proof}

\subsection{Uygulama: Taban Dönüşüm Kodu}

Aşağıdaki C programı, pozitif bir onluk tamsayının $b$ ($2 \le b \le 16$) tabanındaki gösterimini ardışık bölme algoritmasıyla hesaplar:

\begin{lstlisting}[language=CTemel]
#include <stdio.h>

void tabana_cevir(int n, int b) {
    char rakamlar[] = "0123456789ABCDEF";
    char sonuc[64];
    int indis = 0;

    if (n == 0) {
        printf("0\n");
        return;
    }

    while (n > 0) {
        sonuc[indis++] = rakamlar[n % b];
        n = n / b;
    }

    // Kalanlar tersten yazdirilmali
    printf("(%d) tabanindaki degeri: ", b);
    for (int i = indis - 1; i >= 0; i--) {
        putchar(sonuc[i]);
    }
    printf("\n");
}
\end{lstlisting}

\begin{example}[En Sık Yapılan Hatalar]
Taban dönüşümlerinde sıklıkla karşılaşılan hatalar şunlardır:
\begin{enumerate}
    \item \textbf{Kalanları düz sırada yazmak:} Ardışık bölmede ilk elde edilen kalan, sayının en küçük basamağıdır ($b^0$). Dolayısıyla dizi ters çevrilerek okunmalıdır.
    \item \textbf{Geçersiz rakam kullanımı:} Taban $b$ iken basamaklarda $b$ veya daha büyük bir rakam bulunamaz (örneğin $5$ tabanında $(152)_5$ tanımsızdır).
    \item \textbf{Sıfırları atlamak:} Bölüm adımlarında kalan $0$ çıktığında bunu diziye eklemeyi unutmak basamak kaymasına yol açar (örneğin $4 = (100)_2$ iken sıfırları yazmayıp $1$ demek gibi).
\end{enumerate}
\end{example}

\section{Tabanlarda Aritmetik ve İkili Sistemin Rolü}

\subsection{Farklı Tabanlarda Dört İşlem}

Pozisyonel sayı sistemlerinde aritmetik algoritmalar tabandan bağımsızdır: Onluk sistemdeki alt alta toplama, çıkarma ve çarpma kuralları, basamak ağırlığı $10$ yerine $b$ alınarak doğrudan uygulanır.

Toplama ve Elde (Carry) Mekanizması:
Aynı basamaktaki iki rakam toplandığında elde edilen değer $S = a_i + c_i + \text{elde}$ olsun:
\begin{itemize}
    \item Eğer $S < b$ ise basamağa $S$ yazılır, yeni elde $0$ olur.
    \item Eğer $S \ge b$ ise basamağa $S \pmod b$ yazılır ve bir sonraki basamağa $\lfloor S / b \rfloor$ eldesi aktarılır.
\end{itemize}

\begin{example}
$(564)_7$ ile $(345)_7$ sayılarını $7$ tabanında toplayınız.
\end{example}

\begin{proof}[Çözüm]
İşlemi sağdan sola basamak basamak yürütelim:
\begin{itemize}
    \item \textbf{$7^0$ basamağı:} $4 + 5 = 9$. Taban $7$ olduğundan: $9 = 1 \cdot 7 + 2$. Basamağa $2$ yazarız, elde $= 1$.
    \item \textbf{$7^1$ basamağı:} $6 + 4 + 1 \text{ (elde)} = 11$. $11 = 1 \cdot 7 + 4$. Basamağa $4$ yazarız, elde $= 1$.
    \item \textbf{$7^2$ basamağı:} $5 + 3 + 1 \text{ (elde)} = 9$. $9 = 1 \cdot 7 + 2$. Basamağa $2$ yazarız, elde $= 1$.
    \item \textbf{$7^3$ basamağı:} Yalnızca elde kalmıştır: $1$.
\end{itemize}
Sonuç: $(1242)_7$.
\end{proof}

Çıkarma ve Borç (Borrow) Mekanizması:
Bir basamaktaki eksilen rakam çıkan rakamdan küçükse, sol basamaktan ``$1$ borç'' alınır. Onluk tabanda borç alındığında basamağa $10$ eklenirken, $b$ tabanında o basamağa \textbf{$b$} eklenir.

\begin{example}
$(1001)_2 - (0111)_2$ çıkarma işlemini ikili tabanda yapınız.
\end{example}

\begin{proof}[Çözüm]
Sağdan sola inceleyelim:
\begin{itemize}
    \item \textbf{$2^0$ basamağı:} $1 - 1 = 0$.
    \item \textbf{$2^1$ basamağı:} $0 - 1$ yapılamaz, soldan borç aranır. $2^2$ basamağında $0$ olduğundan borç $2^3$ basamağındaki $1$'den alınır ve sağa aktarılır. $2^1$ basamağına $2$ eklenir: $2 - 1 = 1$.
    \item \textbf{$2^2$ basamağı:} Borç aktarımı sonrasında burada $1$ kalmıştır: $1 - 1 = 0$.
    \item \textbf{$2^3$ basamağı:} $1$ borç verildiğinden geriye $0$ kalmıştır: $0 - 0 = 0$.
\end{itemize}
Sonuç: $(0010)_2 = (10)_2$.
Onluk sistemde sağlama: $9 - 7 = 2$.
\end{proof}

\subsection{Genel Tabanlarda Bölünebilme Kuralları}

Onluk sistemden bildiğimiz $9$'a bölünebilme kuralı (rakamlar toplamının $9$'a bölünmesi), $10 \equiv 1 \pmod 9$ bağıntısının bir sonucudur. Bölüm 14'te ele aldığımız modüler aritmetik araçlarını taban aritmetiğiyle birleştirdiğimizde genel bölünebilme kuralları kendiliğinden ortaya çıkar.

\begin{theorem}[$b-1$ ve $b+1$ ile Bölünebilme]
$N = (a_k a_{k-1} \dots a_1 a_0)_b$ sayısı verilsin.
\begin{enumerate}
    \item $N \equiv \sum_{i=0}^k a_i \pmod{b-1}$. Yani $N$, basamakları toplamı ile $(b-1)$ modülünde denktir.
    \item $N \equiv \sum_{i=0}^k (-1)^i a_i \pmod{b+1}$. Yani $N$, birler basamağından başlayarak bir artı bir eksi alınarak kurulan almaşık basamak toplamı ile $(b+1)$ modülünde denktir.
\end{enumerate}
\end{theorem}

\begin{proof}
$b \equiv 1 \pmod{b-1}$ olduğundan her $i \ge 0$ için $b^i \equiv 1^i \equiv 1 \pmod{b-1}$ olur. Polinom açılımında her $b^i$ yerine $1$ yazarsak:
\[
N = \sum_{i=0}^k a_i b^i \equiv \sum_{i=0}^k a_i (1)^i = \sum_{i=0}^k a_i \pmod{b-1}
\]
İkinci kısım için, $b \equiv -1 \pmod{b+1}$ olduğundan $b^i \equiv (-1)^i \pmod{b+1}$ elde edilir:
\[
N = \sum_{i=0}^k a_i b^i \equiv \sum_{i=0}^k a_i (-1)^i = a_0 - a_1 + a_2 - a_3 + \dots \pmod{b+1}
\]
İspat tamamlanır.
\end{proof}

Bu teorem pratik sonuçlar üretir:
\begin{itemize}
    \item $10$ tabanında: $10-1 = 9$ için basamaklar toplamı, $10+1 = 11$ için almaşık toplam kuralı çıkar.
    \item $2$ tabanında: $b-1 = 1$ aşikardır; ancak $b+1 = 3$ için kullanışlı bir kural verir: İkili tabandaki bir sayının çift sıradaki bitleri toplamı ile tek sıradaki bitleri toplamının farkı $3$'e bölünüyorsa, sayı $3$'e tam bölünür.
\end{itemize}

\begin{example}
İkili tabanda yazılmış $N = (101011)_2$ sayısı $3$'e tam bölünür mü?
\end{example}

\begin{proof}[Çözüm]
Bitleri sağdan sola ($0$. bitten başlayarak) numaralandıralım:
\begin{itemize}
    \item $a_0 = 1$ (çift)
    \item $a_1 = 1$ (tek)
    \item $a_2 = 0$ (çift)
    \item $a_3 = 1$ (tek)
    \item $a_4 = 0$ (çift)
    \item $a_5 = 1$ (tek)
\end{itemize}

Almaşık toplam:
\[
(a_0 + a_2 + a_4) - (a_1 + a_3 + a_5) = (1 + 0 + 0) - (1 + 1 + 1) = 1 - 3 = -2
\]
$-2 \not\equiv 0 \pmod 3$ olduğundan sayı $3$'e tam bölünmez; kalan $-2 \equiv 1 \pmod 3$'tür.
Onluk sistemde kontrol edelim: $(101011)_2 = 32 + 8 + 2 + 1 = 43$. Gerçekten $43 = 3 \cdot 14 + 1$ olup kalan $1$'dir.
\end{proof}

\subsection{İkili Sistemin Bilgisayar Bilimindeki Rolü}

Bilgisayar biliminde algoritmik verimliliğin temelinde ikili sistem yer alır. Donanım düzeyinde çarpma ve bölme işlemleri görece maliyetliyken, bit kaydırma ve mantıksal işlemler tek bir saat çevriminde ($O(1)$) çalışır.

Bit Kaydırma (Bit Shifting):
\begin{itemize}
    \item \textbf{Sola kaydırma (\texttt{x << k}):} Sayının ikili gösteriminin sağına $k$ tane sıfır ekler. Bu işlem sayıyı $2^k$ ile çarpmaya denktir:
    \[
    x \ll k = x \cdot 2^k
    \]
    \item \textbf{Sağa kaydırma (\texttt{x >> k}):} Sayının son $k$ bitini atar. Bu işlem pozitif tamsayılarda $\lfloor x / 2^k \rfloor$ bölmesine denktir:
    \[
    x \gg k = \lfloor x / 2^k \rfloor
    \]
\end{itemize}

Bit Düzeyinde İşlemler (Bitwise Operators):
İki sayının ikili gösterimlerindeki karşılıklı bitler mantıksal kapılardan geçirilir:
\begin{enumerate}
    \item \textbf{AND (\texttt{\&}):} Yalnızca her iki bit de $1$ ise $1$ üretir.
    \item \textbf{OR (\texttt{|}):} Bitlerden en az biri $1$ ise $1$ üretir.
    \item \textbf{XOR (\texttt{\textasciicircum}):} Bitler birbirinden farklı ise $1$, aynı ise $0$ üretir (eldesiz ikili toplama).
    \item \textbf{NOT (\texttt{\textasciitilde}):} Bitleri tersine çevirir ($0 \to 1, 1 \to 0$).
\end{enumerate}

Kümelerin Bitmask ile Gösterimi:
$n$ elemanlı bir evrensel kümenin $S = \{e_0, e_1, \dots, e_{n-1}\}$ herhangi bir alt kümesi, $n$ bitlik bir ikili sayı ile birebir eşlenebilir:
$i$. eleman kümede varsa $i$. bit $1$, yoksa $0$ yapılır.
Örneğin $n=4$ için $\{e_0, e_2\}$ alt kümesi $(0101)_2 = 5$ tamsayısı ile temsil edilir.
Böylece küme işlemleri doğrudan bit operatörlerine dönüşür:
\begin{itemize}
    \item Kesişim işlemi: \texttt{A \& B}
    \item Birleşim işlemi: \texttt{A | B}
    \item Simetrik fark: \texttt{A \textasciicircum{} B}
    \item Boş küme: \texttt{0}
    \item Evrensel küme: \texttt{(1 << n) - 1}
\end{itemize}

Bu sayede $O(n)$ sürecek küme işlemleri tek bir makine komutuyla $O(1)$ sürede tamamlanır.

\subsection{Hızlı Üs Alma (Binary Exponentiation)}

İkili tabanın algoritmik problem çözmedeki en etkili uygulamalarından biri \textbf{Hızlı Üs Alma} algoritmasıdır. $a^n$ değerini ardışık çarpma ile hesaplamak $n-1$ çarpma, yani $O(n)$ zaman gerektirir. $n = 10^{18}$ gibi büyük girdiler için bu yaklaşım çalışmaz.

$n$ sayısını ikili tabanda açarsak:
\[
n = \sum_{i=0}^k c_i 2^i \implies a^n = a^{\sum c_i 2^i} = \prod_{c_i = 1} a^{2^i}
\]
Her bir $a^{2^i}$ terimi, bir önceki terimin karesi alınarak ($a^{2^{i}} = (a^{2^{i-1}})^2$) tek bir çarpmayla elde edilebilir. Böylece üs alma işlemi $O(\log n)$ adıma iner.

\begin{example}
$3^{13}$ sayısını hızlı üs alma yöntemiyle hesaplayınız.
\end{example}

\begin{proof}[Çözüm]
Üssü ikili tabanda yazalım: $13 = 8 + 4 + 1 = (1101)_2$. Buradan:
\[
3^{13} = 3^8 \cdot 3^4 \cdot 3^1
\]
Ardışık kareleri hesaplayalım:
\begin{align*}
3^1 &= 3 \\
3^2 &= 3^2 = 9 \\
3^4 &= 9^2 = 81 \\
3^8 &= 81^2 = 6561
\end{align*}
İkili gösterimde $1$ olan basamaklara karşılık gelen değerleri çarparsak:
\[
3^{13} = 3^8 \cdot 3^4 \cdot 3^1 = 6561 \cdot 81 \cdot 3 = 6561 \cdot 243 = 1\,594\,323
\]
Normalde $12$ çarpma gerekecekken, $3$ kare alma ve $2$ çarpma ile toplam $5$ işlemde sonuca ulaşılır.
\end{proof}

Aşağıdaki sözde kod, hızlı üs almanın ikili bit kaydırmayla nasıl gerçeklendiğini özetler:

\begin{lstlisting}
FUNCTION modularExponentiation(base, exponent, modulus)
    result = 1
    base = base MOD modulus

    WHILE exponent > 0 DO
        IF (exponent AND 1) == 1 THEN
            result = (result * base) MOD modulus
        END IF
        base = (base * base) MOD modulus
        exponent = exponent >> 1
    END WHILE

    RETURN result
END FUNCTION\end{lstlisting}

\subsection{Kombinatorik ve Sayı Teorisinde İkili Taban: Kummer ve Lucas Sezgisi}

İkili taban, binom katsayılarının tekliği-çiftliği (paritesi) konusunda pratik sonuçlar sunar.

\begin{theorem}[Sierpiński / Lucas Parite Teoremi]
$\binom{n}{k}$ binom katsayısının tek sayı olması için gerek ve yeter koşul, $k$'nin ikili gösterimindeki her bir $1$ bitinin, $n$'nin ikili gösteriminde de aynı pozisyonda $1$ olmasıdır (yani $k \text{ AND } n = k$; başka bir deyişle $k$, bitwise $n$'nin alt kümesi olmalıdır).
\end{theorem}

\begin{proof}[Fikir / Sezgi]
Bölüm 12'de gördüğümüz Legendre formülüne göre bir $m!$ içindeki $p$ asalının kuvveti $v_p(m!) = \sum_{j=1}^\infty \lfloor \frac{m}{p^j} \rfloor$ ile verilir.
Özel olarak $p=2$ için bu toplam:
\[
v_2(m!) = m - S_2(m)
\]
biçimine dönüşür; burada $S_2(m)$, $m$ sayısının ikili tabandaki rakamları toplamıdır ($1$ bitlerinin sayısı, diğer adıyla popcount).

Binom katsayısı $\binom{n}{k} = \frac{n!}{k!(n-k)!}$ olduğundan, $2$ ile bölünebilme derecesi:
\[
v_2\left(\binom{n}{k}\right) = v_2(n!) - v_2(k!) - v_2((n-k)!)
\]
\[
= \big(n - S_2(n)\big) - \big(k - S_2(k)\big) - \big(n-k - S_2(n-k)\big)
\]
\[
= S_2(k) + S_2(n-k) - S_2(n)
\]
olur. İkili tabanda $k$ ile $n-k$'yi toplarken her elde oluştuğunda, iki tane $1$ biti birleşip üst basamağa tek bir $1$ olarak geçer; yani toplam bit sayısı $1$ azalır.
Dolayısıyla $S_2(k) + S_2(n-k) - S_2(n)$, $k$ ile $n-k$ toplanırken oluşan \textbf{toplam elde sayısıdır} (Kummer Teoremi).

$\binom{n}{k}$'nın tek olması için $v_2 = 0$ olmalıdır; bu da toplama sırasında \textbf{hiç elde oluşmaması} anlamına gelir. İkili toplamda hiç elde çıkmaması ise ancak ve ancak $k$'nin $1$ olduğu hiçbir basamakta $n-k$'nin $1$ olmamasıyla mümkündür. Bu da $n$'in o basamakta zorunlu olarak $1$ olmasını gerektirir.
\end{proof}

\begin{example}
$\binom{11}{k}$ sayısının tek sayı olduğu kaç farklı $k \in \{0, 1, \dots, 11\}$ değeri vardır?
\end{example}

\begin{proof}[Çözüm]
Teorem gereği $k$'nin ikili gösterimi, $11$'in ikili gösteriminin bir alt kümesi olmalıdır.
$11$'i ikili tabana çevirelim:
\[
11 = 8 + 2 + 1 = (1011)_2
\]
$11$'in ikili yazımında $3$ tane $1$ biti vardır ($2^3, 2^1, 2^0$ basamakları).
$k$'nin ikili yazımındaki $1$ bitleri yalnızca bu üç pozisyondan seçilebilir; diğer basamaklar ($2^2$ ve daha büyükler) $0$ kalmalıdır.

Bu üç pozisyonun her biri için iki seçenek vardır: $1$ bitini alabiliriz veya almayız.
Dolayısıyla yazılabilecek geçerli $k$ sayısı:
\[
2^3 = 8
\]
farklı değerdir. Bu değerler sırasıyla:
\begin{align*}
(0000)_2 &= 0 \\
(0001)_2 &= 1 \\
(0010)_2 &= 2 \\
(0011)_2 &= 3 \\
(1000)_2 &= 8 \\
(1001)_2 &= 9 \\
(1010)_2 &= 10 \\
(1011)_2 &= 11
\end{align*}
Böylece Pascal üçgeninin $11$. satırında tam olarak $8$ adet tek sayı, $4$ adet çift sayı bulunduğu doğrulanır.
\end{proof}

\subsection{Negatif Tabanlar ve Alternatif Sistemler}

Yarışmalarda zaman zaman tabanın pozitif tamsayı olmadığı durumlar da karşımıza çıkabilir. Öne çıkan örneklerden biri \textbf{negatif tabanlar} ($b = -2$, negabinary) sistemidir. Negatif tabanların dikkat çekici özelliği, \textbf{eksi işaretine gerek kalmadan} hem pozitif hem de negatif tüm tamsayıları tek bir biçimde temsil edebilmesidir.

$b = -2$ olduğunda basamak değerleri:
\[
(-2)^0 = 1, \quad (-2)^1 = -2, \quad (-2)^2 = 4, \quad (-2)^3 = -8, \quad (-2)^4 = 16, \dots
\]
olur. Kalanlar negatif olmayan $\{0, 1\}$ kümesinden seçilmelidir:
\[
N = q \cdot (-2) + r, \quad r \in \{0, 1\}
\]
Standart tamsayı bölmesinde kalan negatif çıkarsa (örneğin $-5 \div (-2) = 2$, kalan $-1$), bölüm $1$ artırılır ve kalan pozitif yapılır:
\[
-5 = 3 \cdot (-2) + 1 \implies q = 3, r = 1
\]

\begin{example}
$-6$ sayısını $-2$ tabanında (negabinary) yazınız.
\end{example}

\begin{proof}[Çözüm]
Kalanların daima $0$ veya $1$ olmasına dikkat ederek ardışık bölme uygulayalım:
\begin{align*}
-6 &= 3 \cdot (-2) + 0 \quad (\text{kalan } 0) \\
3 &= (-1) \cdot (-2) + 1 \quad (\text{kalan } 1) \\
-1 &= 1 \cdot (-2) + 1 \quad (\text{kalan } 1) \\
1 &= 0 \cdot (-2) + 1 \quad (\text{kalan } 1)
\end{align*}
Kalanları aşağıdan yukarıya doğru okursak:
\[
-6 = (1110)_{-2}
\]
Sağlamasını yapalım:
\[
1 \cdot (-2)^3 + 1 \cdot (-2)^2 + 1 \cdot (-2)^1 + 0 \cdot (-2)^0 = -8 + 4 - 2 + 0 = -6
\]
Sonuç doğrulanır.
\end{proof}

Taban aritmetiği yalnızca sayıların yazım biçimini belirlemekle kalmaz; ardışık bölme algoritmasından modüler aritmetiğe, bitwise işlemlerden dinamik programlamanın bitmask optimizasyonlarına kadar algoritmik düşüncenin temel yapı taşlarından birini oluşturur.
